\section{Discussion}
\label{sec:discussion}

\subsection{What actually provides the safety}
\label{sec:disc-safety}
Across every experiment, deterministic mechanisms decide whether an unsafe action
reaches the feeder: the Evidence Gate decides whether anything may act, the
class-aware removals decide what may act, and the never-auto-execute rule
keeps a model-proposed irreversible primitive from executing autonomously. The ablation
localises containment almost entirely in the class-aware avoidance; components
that improve the \emph{model's} behaviour leave the \emph{system's} safety
unchanged, because safety was never delegated to the model. The trust boundary is
placed so the model's error rate multiplies into operator workload and advisory
quality, not physical consequence.

\subsection{What the LLM contributes --- and does not}
\label{sec:disc-llm}
The evaluated LLM advisers do not improve physical outcomes. The projection arms
show no statistically detectable difference from the deterministic fast path on
ENS and curtailment ($+0.19$ and $+0.40$\,kWh versus D0, both Holm $p=1.0$); of
the class-override arms, ClassOnly-L is worse after correction (Holm $p=0.041$)
while ClassOnly-Q is worse in point estimate but not significant (Holm
$p=0.107$); and the advisory path alone misses attacks the deterministic path
catches --- its in-loop classification is far below the deterministic evidence
mapping.
Crucially, the ground-truth-class oracle through the same shield \emph{is}
modestly better than the deterministic path ($-0.29$\,kWh, 95\% CI
$[-0.69, -0.05]$, excluding zero; raw $p=0.015$ but Holm-adjusted
$p=0.11$, so not significant after multiplicity correction). Two things
follow. There is a small amount of headroom that a \emph{correct}
classifier could capture, concentrated in the DoS cases where the
deterministic EH ranking already recovers most of it; and the evaluated
LLMs do not reach even that modest headroom. The result is not that
advisory quality is irrelevant, but that runtime projection contains
unsafe execution while leaving the model's contribution to physical
performance small and, for the models we tested, unrealised. What the model does
contribute is structured hypotheses, rationales for escalation, and an interface
for operator interrogation. The present evaluation establishes \emph{containment}
of the advisory model, not its operational utility: these outputs may support
operator-facing workflows, but their effect on operator accuracy, workload,
response time, and trust was not evaluated, and our groundedness audit
(70--73\% templated rationales, 23--30\% asserting evidence absent from the
input) shows the rationales are imperfect. Those benefits therefore remain
hypotheses for future human-subject study. The defensible conclusion is that
consultation can be isolated from direct execution authority within the tested
architecture; whether that consultation is operationally worthwhile remains
unestablished.

\subsection{Why model output must gate neither action nor inaction}
\label{sec:disc-confidence}
The two sharpest negative results are mirror images. A shield variant that
auto-executed an irreversible action on the model's stated confidence was worse
than no shield, because injection payloads manufacture exactly the
high-confidence prediction the gate keyed on; and a projection that honoured the
model's \textsc{no\_op} stood down a deterministically evidenced DoS, because the
attack shows the model a frozen copy of \emph{normal} telemetry --- the most
plausible input for ``no incident.'' In both directions the attacker does not
defeat the shield; it feeds the model. The rule is symmetric and, we believe,
load-bearing for any LLM-in-the-loop runtime assurance: \emph{no
attacker-influenceable quantity may relax a constraint on an irreversible action
or stand down a deterministically evidenced incident}. Confidence, rationale,
memory, predicted class, and proposed inaction all fail this test; protocol-level
evidence computed outside the model passes it.

\subsection{What the estimator episode teaches}
\label{sec:disc-estimator}
Counterfactual validation showed a plausible 60-second Energy Impact estimator
selected a tie-optimal action in only about a third of configurations, because it
was calibrated over a window on which the five primitives are nearly
indistinguishable while their true costs diverge over the incident horizon. The
repair needed no learning: the same physical priors, integrated over a
dev-calibrated estimate of the \emph{remaining incident duration} with
restoration-horizon accounting, rank optimally on the (small) held-out set. The
lesson is that impact estimators in response systems must be validated against
branched counterfactual rollouts \emph{at the horizon on which their decisions
bind}. A second point: decision-time identifiability limits matter less than they
appear --- replay and DoS cannot be distinguished when the whole telemetry
sample is stale, but the minimax action (\textsc{freeze\_setpoint}) is
optimal under one hypothesis and free under the other, so the ambiguity
need not be resolved to act well.

\subsection{What the evidence-provenance audit taught us}
\label{sec:disc-provenance}
A methodological point underlies the FDI result and generalises. A
tempting way to detect false-data injection in a simulator is to consult a
per-event flag that the attack generator itself sets; such a flag is
categorical and perfectly reliable, and it makes a gate look stronger than
any detector limited to observable evidence could be. We forbid it: every gate feature is
computed only from wire-observable telemetry (Table~\ref{tab:provenance}).
Under that discipline three families keep genuine protocol signatures
(command packets, sequence regression, timestamp staleness), but a
within-capacity FDI --- a reading that is wrong yet not physically
impossible --- has no observable hard signature, and we report FDI coverage
by magnitude and name redundant sensing as the missing ingredient. The
general lesson is that in a simulation-based security evaluation the
decisive audit is of \emph{where every runtime feature comes from}: a
feature sourced from generator ground truth can make an architecture appear
to solve a problem it has not. The runtime path uses no simulator attack labels
or generator ground truth; its wire-observable evidence nonetheless remains
forgeable wherever authentication or redundant sensing is absent, so we do not
claim the detector is as strong as its clean-signal coverage suggests.

\subsection{Can this run under realistic OT timing?}
\label{sec:disc-timing}
Yes, with a specific division of labour. The deterministic path ---
FeatureView, gate, shield --- adds sub-millisecond per-tick cost and
responds within one control tick of sustained evidence
(${\approx}5$\,s with the frozen persistence thresholds). Model
consultation costs ${\approx}5.5$\,s ($K{=}1$) on a consumer GPU, so it
lands within a 10\,s SLO and is skipped under tighter deadlines with no
loss of containment; $K{=}3$ ($\approx$17\,s) is not the default because
it misses even a 15\,s SLO. The gate also rations the model's exposure: relative to the ungated model-only
arms, benign-day consultations fall by more than 80\% (a descriptive comparison
of consultation counts on the 16 benign holdout configurations), saving inference
budget and reducing the model's exposure to attacker-controllable prompt content.
Under DoS the response is honest rather than flattering: the ungated pipeline's
apparent immunity in that family came from benign-day false freezes that happened
to pre-mitigate later attacks; the gated system pays a small measured ENS cost
for not false-acting
on every benign day.

\subsection{Limitations}
\label{sec:disc-limitations}
\begin{enumerate}
\item \textbf{Benchmark and simulation scope.} All cyber-physical
  results come from a StubFeeder co-simulation and static OpenDSS
  snapshot solves on IEEE 13/34-node feeders; no dynamic time-series
  study, no real feeder, no utility deployment. The StubFeeder's
  physics (nominal-dispatch rule, ENS threshold) is exactly the
  structure EH exploits; real-feeder EH would inherit real model error.
\item \textbf{Model and data scope.} One adapter pair (Qwen2.5-7B,
  Llama-3.1-8B) fine-tuned on one 200-example synthetic corpus;
  evaluation prompts are scenario-derived rather than human-annotated;
  results may not transfer to other backbones or corpora.
\item \textbf{Detection-layer scope, and the FDI observability limit.}
  The gate's observable protocol signatures (forged command packets, sequence
  regression, staleness, rating-plausibility residual) presuppose an
  instrumented protocol layer, and --- absent authentication --- are forgeable
  by a network attacker rather than hard proof of an attack. Under the observable-only feature set, a
  \emph{within-capacity} false-data-injection attack has no observable
  hard signature and is not separable from benign operation without a
  redundant measurement channel, which this simulator does not model;
  such FDI is missed (gate-open $0/1$ at medium and low magnitude).
  Benign protocol quirks (historian echo) remain unresolvable by the gate
  alone, and a stealthy attacker who forges \emph{internally consistent}
  protocol evidence is outside the evaluated threat suite.
\item \textbf{Empirical, not formal, assurance.} Zero-execution results
  are suite-bounded (95\% Clopper--Pearson upper bound $0.051$ on 70
  cases); the invariants are deterministic and directly testable but not
  formally proved, and an adaptive white-box adversary is not modelled.
\item \textbf{Human factors.} HITL operators are deterministic
  stand-ins. The observed escalation burden was strongly
  backbone-dependent (SafeProj-Q escalated 6 of 33 attack configurations,
  SafeProj-L none, because Qwen proposes the irreversible primitive far more
  often than Llama), but the available evaluation was not designed to establish
  the cause or the operational acceptability of that difference, and says
  nothing about real operator workload, fatigue, or trust.
\item \textbf{Auditability scope.} The audit chain provides append-only
  integrity verified by tamper tests, not cryptographic timestamping or
  a signature scheme.
\item \textbf{Drift and lifecycle.} Adapters, thresholds and the frozen
  gate calibration are point-in-time; concept drift, fleet
  reconfiguration, and adversary adaptation over time are unevaluated.
\end{enumerate}
